
\section{Related Work}

\looseness=-1
\textbf{Autoregressive audio language models.} Early autoregressive audio models operated on raw waveforms such as WaveNet \citep{wavenet}, learned discrete codes via VQ-VAE, such as in Jukebox \citep{jukebox} or continuous word-sized audio tokens \citep{algayres2023generativespokenlanguagemodel}. The advent of neural audio codecs \citep{soundstream, encodec} enabled high-quality vector-quantized tokenizers for general audio. These, in turn, powered discrete-token audio LMs: AudioLM \citep{audiolm} for unconditioned audio generation while AudioGen~\citep{audiogen}, MusicLM \citep{musiclm}, and MusicGen \citep{musicgen} apply similar methods for text-to-audio and text-to-music generation. In speech, autoregressive modeling of RVQ codes has been used for text-to-speech generation \citep{valle, speartts} as well as for spoken dialogue \citep{moshi} or translation \citep{hibiki}. However, all of these systems rely on lossy quantization, which inevitably degrades audio fidelity unless a significant compute budget is spent to generate a deep hierarchy of RVQ tokens. This is unlike \ours{}, which predicts continuous embeddings in one pass, providing a better quality-computation trade-off.

\looseness=-1
\textbf{Continuous autoregressive models.} In GIVT, \citet{givt} use a transformer to autoregressively model the latent space of a VAE trained on images, parameterize a Gaussian Mixture Model, and train it with a cross-entropy loss. The authors of MAR \citep{autoreg_image_quant} obtain better results by replacing the Gaussian mixture model with a diffusion-model enabling the approximation of more diverse distributions. In MAR, a large transformer backbone predicts a continuous embedding $\mathbf{z}^s$ given $(\mathbf{x}^1, \ldots, \mathbf{x}^{s-1})$, which then conditions a small MLP diffusion network that models the probability distribution $p(\mathbf{x}^{s}|\mathbf{z}^s)$ of the next latent. 
This eliminates the need for discrete tokenizers, but at the cost of slow sampling: MAR typically needs hundreds of denoising steps per token. \citet{fast_MAR} aims to speed up MAR by replacing the diffusion head with a shortcut head \citep{shortcut_model} for few-step sampling. Shortcut models combine a diffusion loss and a self-consistency loss in order to accelerate the diffusion process. They reduce the number of diffusion steps from 100 to 8 with a similar image quality. Remarkably, \ours{} achieves a quality comparable to the best discrete models with only one step of consistency modeling.

\looseness=-1
In audio, the approach of MAR has been adapted for the task of Text-to-Speech (TTS) synthesis. SALAD \citep{salad} introduces a zero-shot TTS model that operates on continuous speech representations using a per-token latent diffusion process. By leveraging semantic tokens for contextual information and determining synthesis stopping points, SALAD achieves improved intelligibility and audio quality without relying on quantization. Similarly, DiTAR \citep{ditar} presents a patch-based autoregressive framework combining a language model with a diffusion transformer for speech generation. This approach models dependencies between aggregated local patches of continuous tokens, using a causal language model to produce embeddings, which, along with previous patches, serve as inputs to a bidirectional diffusion transformer that predicts the next patch. The authors observe that providing local context through patching was determinant to improve their model, which corroborates what we observe with the introduction of a short context transformer into our model. 
In \citet{sony_trick}, the authors apply the MAR framework to music generation using a relatively small dataset comprising 20,000 single-instrument stems, training their model on 10-second excerpts on top of their continuous compression model Music2Latent \citep{music2latent}. They introduce a method for noise augmentation of the data that allows the model to avoid error accumulation. However, we notice that scaling to a more complex and diverse dataset consisting of full musical pieces makes their approach struggle to maintain high-quality generation on longer sequences. To address these challenges, we propose novel strategies that enhance generation quality while also improving inference efficiency. Finally, Music2Latent2~\citep{music2latent2} explores combining autoregressive and consistency modeling but for compression. IMPACT \citep{impact} explores MAR decoding for text-to-audio generation and sets a new standard for short latency models on the AudioCaps \citep{audiocaps} benchmark. More recently, MingUni-Audio \citep{ming_uni_audio} shows that continuous speech language models can scale to 20B Mixture of Experts models with 3B active parameters.

Some other autoregressive speech models work in the continuous domain thanks to spectral representations such as MELLE \citep{melle} for TTS and Flow-Omni \citep{flow_omni} for speech-to-speech conversation. 


% \paragraph{Other continuous autoregressive audio modeling methods}
% MELLE \cite{melle}, KALL-E FELL-E?